\section{Related Work}
\label{sec:related_work}

The literature relevant to our work lies at the intersection of model merging, test-time adaptation, and stochastic optimization on non-convex landscapes.

\subsection{Training-Free Model Merging}
Model merging has gained substantial popularity as a way to combine specialized capabilities from distinct pre-trained checkpoints without the heavy computational cost of multi-task fine-tuning.
The foundational baseline, \emph{Task Arithmetic}~\cite{ilharco2022editing}, adds fine-tuning task vectors (obtained by subtracting the pre-trained weights from fine-tuned weights) to the pre-trained base model.
While effective for orthogonal tasks, simple addition causes severe parameter interference when tasks share highly correlated features or have overlapping pathways.

To mitigate interference, several advanced training-free methods have been proposed.
\emph{Ties-Merging}~\cite{yadav2023ties} resolves conflicts by retaining only the top-$k$\% parameter changes based on magnitude, resolving sign disagreements via majority voting, and scaling the remaining changes.
\emph{DARE}~\cite{jin2023dataless} randomly drops small parameter changes and rescales the remaining ones to preserve the expectation of task vectors.
\emph{ZipIt!}~\cite{zhang2024zipit} and \emph{Git Re-Basin}~\cite{ainsworth2023git} align neuron orderings modulo permutation symmetries before merging.
More recently, \emph{OrthoMerge}~\cite{yang2026orthomerge} preserved geometry by performing model merging on the Riemannian manifold of orthogonal groups.
However, all these methods are training-free and data-free, meaning they rely on static heuristics and cannot adjust merging parameters to adapt to target downstream distributions.

\subsection{Test-Time Adaptive Model Merging}
To overcome the static nature of training-free merging, recent works have pioneered \emph{test-time adaptive model merging}, where parameters are adapted on unlabeled target task data during inference.
\emph{AdaMerging}~\cite{yang2024adamerging} optimizes layer-wise merging coefficients by minimizing the prediction entropy of the merged model.
However, entropy minimization is notorious for training instabilities and can cause degenerate predictions (e.g., collapsing all predictions to a single class to yield zero entropy).

To address this, \emph{SyMerge}~\cite{jung2025symerge} introduced an expert-guided self-labeling cross-entropy proxy loss, using the fixed predictions of the unmerged fine-tuned experts as stable teachers.
SyMerge jointly optimizes the layer-wise merging coefficients $\Lambda$ and a task-specific adapter layer $\Theta^{tr}$ at test-time.
Concurrently, \emph{SAIM}~\cite{foret2020sharpness} developed sharpness-aware isotropic merging for continual learning to balance the singular value spectrum.
While test-time adaptation significantly improves multi-task performance, all existing test-time model merging methods rely on standard deterministic optimization (e.g., Adam or SGD).
Under severe task interference, the multi-task proxy loss landscape is highly non-convex, containing severe ripples and sharp sub-optimal local traps.
Because deterministic optimizers have no mechanism to escape energy barriers, they are highly sensitive to initialization, frequently getting trapped in sub-optimal basins.

\paragraph{Alternative Optimization Heuristics.} One might consider alternative deep learning optimization heuristics designed to find flat minima or escape sharp traps, such as Sharpness-Aware Minimization (SAM)~\cite{foret2020sharpness}, Stochastic Weight Averaging (SWA)~\cite{izmailov2018averaging}, or Cosine Annealing schedules~\cite{loshchilov2016sgdr}.
However, we highlight that these methods are fundamentally unsuitable for the unique constraints of test-time model merging.
First, SAM requires a two-step gradient computation per parameter update to compute an adversarial perturbation, which doubles the forward/backward pass latency during test-time adaptation where real-time inference constraints are paramount.
Furthermore, under unlabeled proxy losses, SAM's adversarial perturbations can easily exploit the self-labeled teacher model, leading to adversarial optimization collapse.
Second, SWA averages weights along the trajectory but requires a long training cycle (hundreds of epochs) to accumulate checkpoints inside flat basins; in contrast, test-time adaptation is extremely data-scarce and operates over only a few epochs and batches (e.g., 5 epochs over 8 batches), making trajectory averaging statistically biased and ineffective.
Finally, learning rate decay schemes like Cosine Annealing remain purely deterministic and cannot escape local energy barriers; scaling down the learning rate in a sub-optimal trap simply freezes the model in that trap faster.
This necessitates physical, isotropic noise injection as proposed in ThermoMerge.

\subsection{Stochastic Gradient Langevin Dynamics \& Simulated Annealing}
Stochastic Gradient Langevin Dynamics (SGLD)~\cite{welling2011bayesian} was originally introduced as an efficient Markov Chain Monte Carlo (MCMC) algorithm for Bayesian posterior sampling.
SGLD combines stochastic gradient descent with coordinate-wise Gaussian noise scaled by the learning rate and a temperature parameter.
In recent years, deep learning optimization researchers have leveraged SGLD and Langevin diffusion not only for posterior sampling, but as a powerful mechanism for global optimization on highly non-convex surfaces~\cite{zhang2020sgld,mandt2017stochastic}.

By injecting temperature-scaled stochastic noise, SGLD behaves like a continuous physical diffusion process (the Langevin equation), allowing parameters to hop over high energy barriers and escape sharp local traps.
Theoretical analyses of loss landscapes~\cite{hochreiter1997flat,keskar2016large,dinh2017sharp} have shown that flat minima exhibit superior generalization and out-of-distribution (OOD) robustness.
SGLD naturally biasses optimization trajectories toward wider, flatter valleys because sharp valleys are "hotter" and have higher escaping probability under thermal agitation.

To ensure global exploration followed by precise local convergence, we combine SGLD with \emph{Simulated Annealing}~\cite{kirkpatrick1983optimization,dekkers1991global}, where the temperature exponentially cools over time.
While Simulated Annealing and Langevin diffusion have been applied to standard neural network training, to the best of our knowledge, we are the first to transfer these thermodynamic optimization principles to the unique and highly non-convex landscape of test-time model merging.